\documentclass[a4paper,11pt]{article}
\usepackage[utf8]{inputenc}
\usepackage[T1]{fontenc}
\usepackage[spanish]{babel}
\usepackage{geometry}
\usepackage{array}
\usepackage{colortbl}
\usepackage{xcolor}
\usepackage{booktabs}
\usepackage{enumitem}
\usepackage{fancyhdr}
\usepackage{tikz}

\geometry{left=2cm, right=2cm, top=2.5cm, bottom=2cm}

\definecolor{violeta}{HTML}{7B2D8E}
\definecolor{azulg}{HTML}{3498DB}
\definecolor{rojog}{HTML}{E74C3C}
\definecolor{gris}{HTML}{F2F2F2}
\definecolor{verde}{HTML}{27AE60}
\definecolor{naranja}{HTML}{E67E22}

\pagestyle{fancy}
\fancyhf{}
\fancyhead[L]{\small\textbf{Nuevo Compañeros 1 — U03: La Familia}}
\fancyhead[R]{\small\textbf{Píldora formativa 3.1}}
\fancyfoot[C]{\thepage}

\setlength{\parindent}{0pt}

\begin{document}

\begin{center}
{\LARGE\textbf{PÍLDORA FORMATIVA 3.1}}\\[0.3cm]
{\Large Patrón de género -o/-a en el vocabulario de parentesco}\\[0.2cm]
{\normalsize Sección Vocabulario (p.34--35) — Proyectable — 3 diapositivas}\\[0.2cm]
{\small Versión enriquecida secuencial para Fases 3 y 5}
\end{center}

\vspace{0.8cm}

%% ============================================================
%% DIAPOSITIVA 1
%% ============================================================

\noindent\textcolor{violeta}{\rule{\linewidth}{2.5pt}}

\section*{DIAPOSITIVA 1 — Patrón regular: -o (masculino) / -a (femenino)}

{\small Fase 3 — Active la conciencia gramatical}

\vspace{0.3cm}

La mayoría de los términos de parentesco sigue un patrón regular de género. El masculino termina en \textbf{-o} y el femenino en \textbf{-a}.

\vspace{0.5cm}

\begin{center}
\renewcommand{\arraystretch}{1.6}
\begin{tabular}{|>{\centering\arraybackslash}p{5cm}|>{\centering\arraybackslash}p{5cm}|}
\hline
\rowcolor{azulg!15}
\textbf{\textcolor{azulg}{Masculino (-o)}} & \textbf{\textcolor{rojog}{Femenino (-a)}} \\
\hline
abuel\textbf{\textcolor{azulg}{o}} & abuel\textbf{\textcolor{rojog}{a}} \\
\hline
herman\textbf{\textcolor{azulg}{o}} & herman\textbf{\textcolor{rojog}{a}} \\
\hline
hij\textbf{\textcolor{azulg}{o}} & hij\textbf{\textcolor{rojog}{a}} \\
\hline
tí\textbf{\textcolor{azulg}{o}} & tí\textbf{\textcolor{rojog}{a}} \\
\hline
prim\textbf{\textcolor{azulg}{o}} & prim\textbf{\textcolor{rojog}{a}} \\
\hline
niet\textbf{\textcolor{azulg}{o}} & niet\textbf{\textcolor{rojog}{a}} \\
\hline
sobrin\textbf{\textcolor{azulg}{o}} & sobrin\textbf{\textcolor{rojog}{a}} \\
\hline
\end{tabular}
\end{center}

\vspace{0.5cm}

\begin{center}
\fbox{\parbox{12cm}{\centering\large
\textcolor{azulg}{-o} = masculino \hspace{2cm} \textcolor{rojog}{-a} = femenino\\[0.2cm]
{\normalsize La terminación cambia, la raíz no.}
}}
\end{center}

\vspace{0.3cm}

{\small\textit{Conexión con conocimiento previo:} El estudiante ya conoce el patrón -o/-a de los adjetivos de nacionalidad (U02: español/española, colombiano/colombiana).}

\newpage

%% ============================================================
%% DIAPOSITIVA 2
%% ============================================================

\noindent\textcolor{violeta}{\rule{\linewidth}{2.5pt}}

\section*{DIAPOSITIVA 2 — Excepciones: raíces distintas}

{\small Fase 3 — Active la conciencia gramatical}

\vspace{0.3cm}

No todos los parentescos siguen el patrón -o/-a. Hay dos pares con \textbf{raíces completamente distintas} que hay que memorizar.

\vspace{0.5cm}

\begin{center}
\renewcommand{\arraystretch}{1.8}
\begin{tabular}{|>{\centering\arraybackslash}p{5cm}|>{\centering\arraybackslash}p{5cm}|}
\hline
\rowcolor{azulg!15}
\textbf{\textcolor{azulg}{Masculino}} & \textbf{\textcolor{rojog}{Femenino}} \\
\hline
\rowcolor{naranja!10}
\textbf{\Large padre} & \textbf{\Large madre} \\
{\small (no *padro)} & {\small (no *madra)} \\
\hline
\rowcolor{naranja!10}
\textbf{\Large marido} & \textbf{\Large mujer} \\
{\small (no *marida)} & {\small (= esposa / wife)} \\
\hline
\end{tabular}
\end{center}

\vspace{0.8cm}

\begin{center}
\fbox{\parbox{12cm}{\centering\large
Estas palabras no siguen -o / -a.\\[0.2cm]
{\normalsize Son pares con raíces distintas. Hay que memorizarlos.}
}}
\end{center}

\vspace{0.8cm}

\begin{center}
\renewcommand{\arraystretch}{1.4}
\begin{tabular}{|l|l|}
\hline
\rowcolor{gris}
\textbf{Patrón regular (-o/-a)} & \textbf{Excepción (raíces distintas)} \\
\hline
abuelo / abuela & padre / madre \\
hermano / hermana & marido / mujer \\
hijo / hija & \\
tío / tía & \\
primo / prima & \\
nieto / nieta & \\
sobrino / sobrina & \\
\hline
\end{tabular}
\end{center}

\vspace{0.5cm}

{\small\textit{Nota sobre mujer:} Tiene doble significado en español: mujer = woman (general) y mujer = wife (esposa). En esta sección se usa como esposa.}

\newpage

%% ============================================================
%% DIAPOSITIVA 3
%% ============================================================

\noindent\textcolor{violeta}{\rule{\linewidth}{2.5pt}}

\section*{DIAPOSITIVA 3 — Plurales mixtos, verbo \textit{ser} y doble significado de \textit{mujer}}

{\small Fase 5 — Amplíe el vocabulario con traducción reflexiva}

\vspace{0.5cm}

%% --- PARTE A: Plurales mixtos ---

\subsection*{\textcolor{verde}{A. Plurales mixtos}}

Cuando el grupo incluye masculino y femenino, el español usa el \textbf{masculino plural}.

\vspace{0.3cm}

\begin{center}
\renewcommand{\arraystretch}{1.5}
\begin{tabular}{|>{\centering\arraybackslash}p{3cm}|>{\centering\arraybackslash}p{1cm}|>{\centering\arraybackslash}p{3cm}|>{\centering\arraybackslash}p{1cm}|>{\centering\arraybackslash}p{3cm}|}
\hline
\rowcolor{azulg!10}
\textbf{\textcolor{azulg}{Masculino}} & & \textbf{\textcolor{rojog}{Femenino}} & & \textbf{Plural mixto} \\
\hline
padre & + & madre & = & \textbf{padres} \\
\hline
abuelo & + & abuela & = & \textbf{abuelos} \\
\hline
hijo & + & hija & = & \textbf{hijos} \\
\hline
nieto & + & nieta & = & \textbf{nietos} \\
\hline
hermano & + & hermana & = & \textbf{hermanos} \\
\hline
tío & + & tía & = & \textbf{tíos} \\
\hline
\end{tabular}
\end{center}

\vspace{0.5cm}

%% --- PARTE B: Doble significado de mujer ---

\subsection*{\textcolor{naranja}{B. Doble significado de \textit{mujer}}}

\begin{center}
\renewcommand{\arraystretch}{1.4}
\begin{tabular}{|p{6cm}|p{6cm}|}
\hline
\rowcolor{naranja!10}
\textbf{mujer = woman} (general) & \textbf{mujer = wife} (esposa) \\
\hline
Significado general: una mujer, las mujeres. & En contexto de parentesco: Alicia es la \textbf{mujer} de Luis. \\
\hline
\multicolumn{2}{|c|}{\cellcolor{gris}\small En esta sección se usa con el significado de \textbf{esposa}.} \\
\hline
\end{tabular}
\end{center}

\vspace{0.5cm}

%% --- PARTE C: Verbo ser ---

\subsection*{\textcolor{violeta}{C. El verbo \textit{ser} en las presentaciones de parentesco}}

En las presentaciones de parentesco siempre se usa el verbo \textbf{ser}.

\vspace{0.3cm}

\begin{center}
\renewcommand{\arraystretch}{1.5}
\begin{tabular}{|>{\centering\arraybackslash}p{7cm}|>{\centering\arraybackslash}p{3cm}|}
\hline
\rowcolor{violeta!10}
\textbf{Ejemplo} & \textbf{Forma de \textit{ser}} \\
\hline
Alicia \textbf{es} la mujer de Luis. & \textbf{es} (singular) \\
\hline
Carlos y Juana \textbf{son} los abuelos de David. & \textbf{son} (plural) \\
\hline
David \textbf{es} el hijo de Alicia y Luis. & \textbf{es} (singular) \\
\hline
Álvaro y Paloma \textbf{son} los primos de David. & \textbf{son} (plural) \\
\hline
\end{tabular}
\end{center}

\vspace{0.5cm}

\begin{center}
\fbox{\parbox{12cm}{\centering\large
Parentesco = verbo \textbf{ser}\\[0.2cm]
{\normalsize 1 persona $\rightarrow$ \textbf{es} \hspace{2cm} 2+ personas $\rightarrow$ \textbf{son}}
}}
\end{center}

\end{document}
